\documentclass{article}
\usepackage[utf8]{inputenc}
\usepackage[T1]{fontenc}
\usepackage{geometry}
\geometry{
    a4paper,
    left=2.5cm,
    right=2.5cm,
    top=2.5cm,
    bottom=2.5cm
}
\usepackage{graphicx}
\usepackage{float}
\usepackage{hyperref}
\usepackage[english]{babel}
\usepackage{amsmath}
\usepackage{amssymb}
\usepackage{array}
\usepackage{booktabs}
\usepackage{caption}
\usepackage{subcaption}
\usepackage{enumitem}
\usepackage{fancyhdr}
\usepackage{titlesec}
\usepackage{url}
\usepackage{textcomp}
\usepackage{xcolor}
\usepackage{listings}
\usepackage{multirow}
\usepackage{longtable}
\usepackage{tabularx}
\usepackage{times}

% Set up page style
\pagestyle{fancy}
\fancyhf{}
\rfoot{\thepage}
\renewcommand{\headrulewidth}{0pt}

% Title formatting
\titleformat{\chapter}[display]{\normalfont\bfseries}{\chaptertitle\thechapter}{0pt}{\Huge}
\titlespacing*{\chapter}{0pt}{50pt}{40pt}

% Document info
\title{Template Document}
\author{Academic Support Agent}
\date{\today}

\begin{document}
\maketitle

\begin{center}
\begin{tikzpicture}[remember picture,overlay]
\node[anchor=center, inner sep= 0pt] 
{\includegraphics[width=.5\textwidth]
{logo_centenario-eps-converted-to.pdf}};
\centering
\end{tikzpicture}
\vfill
\LARGE
\textsc{Universit  degli Studi di Trieste}\\
\rule{\textwidth}{0.1mm}\\
\large
\textsc{Dipartimento di Scienze Economiche, Aziendali, Matematiche e Statistiche - DEAMS}\\ 
\bigskip
Corso di Laurea in Quello Che Vi Pare a Voi \\
\vfill %%\vspace{3cm}
\textsc{Tesi di laurea triennale}\\
\vfill %%\vspace{3cm}
\LARGE
\textbf{\textsc{Derivatives Pricing and Strategies:\\A Focus On Agricultural Commodities}}\\
\end{center}
\vfill
%%%%%%
%%%%%% Qui va inserito il nome del laureando
\begin{tabular}[t]{l}
Laureando:\\
\textbf{Pietro Bettini}\\
\textsc{Matricola: EC12062002}
\end{tabular}
%%%%%%
\hfill 
%%%%%% Qui va inserito il nome del relatore
\begin{tabular}[t]{l}
Relatore: \\
\textbf{Prof. Nico Zanusso} \\
\\
Correlatore: \\
\textbf{Prof. Fabrizio Riccio}
\end{tabular}
%%%%%%
\vfill
\begin{center}
\normalsize
\rule{8cm}{0.1mm}\\
\bigskip
ANNO ACCADEMICO 2023--2024
\end{center}

\large



%%%%% DEDICATION %%%%%
\newpage
\thispagestyle{empty}
\begin{center}
    \textit{A Francesca}
\end{center}
\newpage



%%%%% INDEX AND LIST OF FIGURES %%%%%
\pagestyle{plain} 
\pagenumbering{roman}







%%%%% 
\newpage
\pagestyle{fancy}
\pagenumbering{arabic}
 
 
%%%%% Examples of chapters
 
 
\chapter{Introduction}
This chapter is on how to upload an image and refer to it. See figure \ref{fig:GBM}. Of course you gotta add some text otherwise the page will be empty and the figure will be automatically placed at the top or at the bottom of the page.
% This is how to upload an image
\begin{figure}
    \centering
    \includegraphics[scale=0.65]{GBM comparison.pdf}
    \caption{Comparison between low variance (left) and high variance (right) Geometric Brownian Motion.}
    \label{fig:GBM}
\end{figure}

\section{This is how to upload a figure}
\subsection{And refer to a figure}



\chapter{How to create a table}
\section{How to cite a table}
This is how to create a table, and cite a table. See table \ref{tab:table3}. 

\begin{table}
\centering
\scriptsize % remove this if you want a larger table
    \begin{tabular}{@{}lllll@{}}
\hline
Commodity & Product Code & Price Quotation & Contract size &  \\ \hline
Class III Milk & DC/DA & U.S. dollars and cents per cwt & 2,000 cwt &  \\
Class IV Milk & GDK/DK & U.S. dollars and cents per cwt & 2,000 cwt & \\
Non-fat Dry Milk & GNF & U.S. cents per pound & 44,000 lb &  \\ 
Dry Whey & DY & U.S. cents per pound & 44,000 lb & \\
Cash-settled Butter & CB & U.S. cents per pound & 20,000 lb & \\
Cash-Settled Cheese & CSC & U.S. dollars and cents per pound & 20,000 lb & \\ 
\hline
\end{tabular}
\caption {Dairy. \\
Author s elaborations based on data from CME group.}
\label{tab:table3}
\end{table}

\section{Now that you see the table}
\subsection{You can copy it}
Of course you gotta add some text otherwise the page will be empty and the table will be automatically placed at the top or at the bottom of the page.

\chapter{I sette nani}

This chapter is completely useless.
\section{Mammolo Eolo}
\subsection{Pisolo}
\subsubsection{Non me li ricordo}
\subsection{Brontolo}
\section{Stronzolo}
\section{Jeremy Clarkson}
\section{James May}
\section{Richard Hammond}

\chapter{Other Stuff}

\section{How to draw a binomial tree.}
This chapter shows how you can create a recombiant binomial tree. 

\begin{figure}[htbp]
    \centering
    \begin{tikzpicture}[>=stealth,sloped]
        \matrix (tree) [%
          matrix of nodes,
          minimum size=0.3cm,
          column sep=0.3cm,
          row sep=0.3cm,
        ]
        {
              &    &          &          & $\empty$ &\\
              &    &          & $Su^{3}$ & $\empty$ &\\
              &    & $Su^{2}$ &          & $\empty$ &\\
              & $Su$ &          & $Su^{2}d$& $\empty$ &\\    
          S   &    & $Sud$      &          & $\empty$ &...&\\ 
              & $Sd$ &          & $Sud^{2}$& $\empty$ &\\     
              &    & $Sd^{2}$ &          & $\empty$ &\\
              &    &          & $Sd^{3}$ & $\empty$ &\\
              &    &          &          & $\empty$ &\\
        };
        \draw[->] (tree-5-1) -- (tree-4-2) node [midway,above] {\empty};
        \draw[->] (tree-5-1) -- (tree-6-2) node [midway,below] {\empty};
        \draw[->] (tree-4-2) -- (tree-3-3) node [midway,above] {\empty};
        \draw[->] (tree-4-2) -- (tree-5-3) node [midway,below] {\empty};
        \draw[->] (tree-6-2) -- (tree-5-3) node [midway,above] {\empty};
        \draw[->] (tree-6-2) -- (tree-7-3) node [midway,below] {\empty};
        \draw[->] (tree-3-3) -- (tree-2-4) node [midway,above] {\empty};
        \draw[->] (tree-3-3) -- (tree-4-4) node [midway,below] {\empty};
        \draw[->] (tree-5-3) -- (tree-4-4) node [midway,above] {\empty};
        \draw[->] (tree-5-3) -- (tree-6-4) node [midway,below] {\empty};
        \draw[->] (tree-7-3) -- (tree-6-4) node [midway,above] {\empty};
        \draw[->] (tree-7-3) -- (tree-8-4) node [midway,below] {\empty};
        \draw[->] (tree-2-4) -- (tree-1-5) node [midway,above] {\empty};
        \draw[->] (tree-2-4) -- (tree-3-5) node [midway,below] {\empty};
        \draw[->] (tree-4-4) -- (tree-3-5) node [midway,above] {\empty};
        \draw[->] (tree-4-4) -- (tree-5-5) node [midway,below] {\empty};
        \draw[->] (tree-6-4) -- (tree-5-5) node [midway,above] {\empty};
        \draw[->] (tree-6-4) -- (tree-7-5) node [midway,below] {\empty};
        \draw[->] (tree-8-4) -- (tree-7-5) node [midway,above] {\empty};
        \draw[->] (tree-8-4) -- (tree-9-5) node [midway,below] {\empty};
    \end{tikzpicture} 
    \caption{Markov process evolution}
    \label{fig:stochasticprocessevolution}
\end{figure}



\section{How to make a bullet list.}
To better understand the pricing models let's first introduce some statistical concepts. Brownian motion $W_t$ is a stochastic process with four main features: 
\begin{itemize}
    \item Starts at time $t$: $W_0 = 0$;
    \item $W_t$ is continuous in time;
    \item $W_t$ increments (the differences $X(t)-X(s)$ between two intervals in time $t$ and $s$) are independent;
    \item $W_t - W_s \sim \mathcal{N}(0, s-u), (\forall 0 < s < t) $
\end{itemize}




%%%%% aggiungi lessico
\chapter{Glossary}
Here you can write your glossary.


\chapter{Ringraziamenti}

Un grazie che non pu  essere taciuto va a Luigi Magliocchi, faro nel buio di UniTs.
\end{document}